\documentclass[11pt]{article}
\usepackage[T1]{fontenc}
\usepackage[dvipsnames]{xcolor}
\usepackage{amsmath,amsthm,mathtools}
\usepackage{newtxtext,newtxmath}
\usepackage[margin=0.93in,headheight=15pt]{geometry}
\usepackage{microtype}
\usepackage{booktabs,array,longtable}
\usepackage{enumitem}
\usepackage{fancyhdr}
\usepackage[most]{tcolorbox}
\usepackage{listings}
\usepackage{hyperref}
\usepackage[nameinlink,noabbrev]{cleveref}
\definecolor{Ink}{HTML}{18364B}
\definecolor{Accent}{HTML}{176C79}
\definecolor{Pale}{HTML}{F0F6F7}
\hypersetup{colorlinks=true,linkcolor=Ink,citecolor=Accent,urlcolor=Accent,
 pdftitle={The Exponential Scale of Gregory Sign Thresholds},
 pdfauthor={GPT-6 Astra Pro -- research draft},pdfsubject={A counterexample and an all-orders asymptotic expansion}}
\pagestyle{fancy}\fancyhf{}
\fancyhead[L]{\small\scshape Gregory sign thresholds}
\fancyhead[R]{\small Research draft \;|\; 19 September 2026}
\fancyfoot[C]{\thepage}
\renewcommand{\headrulewidth}{0.3pt}
\renewcommand{\footrulewidth}{0pt}
\setlength{\parskip}{3pt}
\setlength{\emergencystretch}{2em}
\numberwithin{equation}{section}
\newtheorem{theorem}{Theorem}[section]
\newtheorem{lemma}[theorem]{Lemma}
\newtheorem{proposition}[theorem]{Proposition}
\newtheorem{corollary}[theorem]{Corollary}
\theoremstyle{definition}
\newtheorem{definition}[theorem]{Definition}
\newtheorem{remark}[theorem]{Remark}
\crefname{lemma}{lemma}{lemmas}
\Crefname{lemma}{Lemma}{Lemmas}
\crefname{proposition}{proposition}{propositions}
\Crefname{proposition}{Proposition}{Propositions}
\crefname{corollary}{corollary}{corollaries}
\Crefname{corollary}{Corollary}{Corollaries}
\crefname{remark}{remark}{remarks}
\Crefname{remark}{Remark}{Remarks}
\newcommand{\Gr}{\operatorname{Gr}}
\newcommand{\E}{\mathbb E}
\newcommand{\Prob}{\mathbb P}
\newcommand{\Q}{\mathbb Q}
\newcommand{\R}{\mathbb R}
\newcommand{\N}{\mathbb N}
\newcommand{\dd}{\,\mathrm dt}
\newcommand{\ii}{\mathrm i}
\newcommand{\fall}[2]{#1^{\underline{#2}}}
\newcommand{\rf}[2]{#1^{\overline{#2}}}
\newcommand{\sgn}{\operatorname{sgn}}
\newcommand{\coef}[2]{[#1^{#2}]}
\newtcolorbox{resultbox}{colback=Pale,colframe=Accent,boxrule=0.6pt,
 arc=1mm,left=10pt,right=10pt,top=8pt,bottom=8pt}
\lstset{basicstyle=\ttfamily\small,breaklines=true,columns=fullflexible,
 frame=single,rulecolor=\color{black!15},backgroundcolor=\color{black!2},
 showstringspaces=false}
\title{\color{Ink}\bfseries The Exponential Scale of\protect\\ Gregory Sign Thresholds\protect\\[5pt]
\large A counterexample to a recent conjecture,\protect\\ with an all-orders asymptotic expansion}
\author{Research article prepared with GPT-6 Astra Pro}
\date{19 September 2026}

\begin{document}
\maketitle
\begin{abstract}
For the higher-order Gregory coefficients
\[
 \left(\frac{x}{\log(1+x)}\right)^m=\sum_{n\ge0}G_n^{(m)}x^n,
\]
let $\Gr(m)$ be the least index after which
$(-1)^{n-1}G_n^{(m)}$ is always strictly positive. A September 2026
conjecture places this threshold between $0.47\,3^m$ and $0.477\,3^m$.
We obtain instead
\[
 \log\Gr(m)=m+1-\gamma-\frac{\zeta(2)+\zeta(3)}{m}+O(m^{-2}),
\]
and give a complete asymptotic expansion with an explicit coefficient
algorithm. In particular, the exponential growth factor is $e$, not $3$.
The proof represents a coefficient as an integral over the density of a
sum of independent uniform random variables. Near the last sign change,
the integral reduces to a gamma expectation concentrated around a simple
zero of $-1/\Gamma(-t)$. A cumulative-positivity lemma then converts a local
asymptotic calculation into positivity of \emph{every} subsequent integer
coefficient. Independent exact rational certificates establish
$\Gr(3)=11$, $\Gr(4)=36$, $\Gr(5)=113$, and $\Gr(6)=346$.
Thus there is also a finite counterexample to the conjecture as printed.
The accompanying code separates exact certification, symbolic algebra,
and non-certified numerical exploration.
\end{abstract}

\begin{resultbox}
\textbf{Principal conclusion.}
Writing $d=1-\gamma$,
\[
 \boxed{\Gr(m)\sim e^d e^m,\qquad e^d=1.52620511159586388\ldots .}
\]
The article supplies a proof, not merely a numerical extrapolation.
It is a research draft: the argument has not been independently refereed
or formalized in a proof assistant. The claimed contribution is the
uniform large-order threshold analysis, not the rediscovery of reported
small values or the elementary identification of a polynomial's roots.
\end{resultbox}

\section{The problem and the result}\label{sec:intro}

\subsection{A precise threshold}
All logarithms are natural logarithms, and $\N=\{1,2,\ldots\}$. For an integer
$m\ge1$, define rational numbers $G_n^{(m)}$ by
\begin{equation}\label{eq:gf}
 F_m(x):=\left(\frac{x}{\log(1+x)}\right)^m
       =\sum_{n=0}^{\infty}G_n^{(m)}x^n,
 \qquad G_0^{(m)}=1.
\end{equation}
The expression at $x=0$ is interpreted by analytic continuation.
Set
\begin{equation}\label{eq:threshold}
 a_m(n):=(-1)^{n-1}G_n^{(m)}\quad(n\ge1),\qquad
 \Gr(m):=\min\{N\in\N:a_m(n)>0\text{ for every integer }n\ge N\}.
\end{equation}
Strict positivity is part of the definition: a zero is an obstruction.
We will prove that this minimum exists, independently of the source
paper's eventual-alternation theorem.

The target is Conjecture 6.1 of Xu and Zhao's preprint
\cite{XZ2026}, version 1, submitted on 10 September 2026. In the notation
of this article it asserts
\begin{equation}\label{eq:conjecture}
 0.47\,3^m<\Gr(m)<0.477\,3^m\qquad(m\ge5).
\end{equation}
Their discussion concerns the \emph{exact} onset of permanent alternation;
we make that minimality explicit in \eqref{eq:threshold}. An arbitrary
nonminimal eventual-positivity bound would not define a unique sequence.

There is an immediate warning sign. The source already reports
$\Gr(5)=113$, whereas $0.47\,3^5=114.21$. This discrepancy could be a
starting-index or transcription error. Merely pointing it out would not
explain the proposed growth law. Our main result shows that no change to
a finite starting range can repair the base $3$.

\subsection{Main theorem}
Write $\gamma$ for Euler's constant, $\zeta$ for the Riemann zeta function,
and $d=1-\gamma$.

\begin{theorem}[Threshold asymptotics]\label{thm:main}
As $m\to\infty$ through positive integers,
\begin{equation}\label{eq:main}
 \log\Gr(m)=m+d-\frac{\zeta(2)+\zeta(3)}{m}
                +\frac{C_2}{m^2}+O(m^{-3}),
\end{equation}
where
\begin{equation}\label{eq:C2}
 C_2=\zeta(2)\zeta(3)+\zeta(2)-\zeta(3)
                 -\frac52\zeta(4)-3\zeta(5).
\end{equation}
More generally there are explicitly computable constants $C_j$ such that,
for each fixed $R\ge0$,
\begin{equation}\label{eq:allorders}
 \log\Gr(m)=m+d+\sum_{j=1}^{R}\frac{C_j}{m^j}
                      +O_R(m^{-R-1}).
\end{equation}
Here $C_1=-\zeta(2)-\zeta(3)$ and
$C_j\in\Q[\zeta(2),\ldots,\zeta(2j+1)]$.
\end{theorem}

The all-orders assertion is an asymptotic expansion: it does not assert
convergence of the infinite series. The constants implicit in the error
terms may depend on the truncation order, but not on $m$.

\begin{corollary}\label[corollary]{cor:base}
The threshold satisfies
\begin{equation}\label{eq:limits}
 \lim_{m\to\infty}\frac{\Gr(m)}{e^m}=e^{1-\gamma},\qquad
 \lim_{m\to\infty}\frac{\Gr(m+1)}{\Gr(m)}=e,\qquad
 \lim_{m\to\infty}\frac{\Gr(m)}{3^m}=0.
\end{equation}
Consequently, the lower inequality in \eqref{eq:conjecture} fails for all
sufficiently large $m$.
\end{corollary}
\begin{proof}
The first two limits follow by exponentiating \eqref{eq:main} and by
taking its difference at consecutive orders. The last follows from
$\Gr(m)/3^m=e^{d+o(1)}(e/3)^m$ and $e<3$.
\end{proof}

\subsection{Proof architecture and status}
The essential difficulty is not finding a plausible zero of a coefficient
integral. A threshold statement quantifies over an infinite tail, so a
local calculation alone is insufficient. We first establish a positivity
propagation criterion. We then prove exponentially small bounds for the
parts of the integral away from its first sign-changing zero. Together,
these tools give a genuine upper bound for \emph{all} later coefficients,
while one negative coefficient supplies the matching lower bound.

The source version and its displayed conjecture were checked on
19 September 2026, including the PDF page containing the statement. A
limited search did not reveal a later resolution. No exhaustive
literature-priority claim is made. In particular, classical gamma-function
identities and elementary properties of uniform sums are used as tools,
not presented as new discoveries. The classical Gregory sequence is
recorded through its numerators and denominators in OEIS entries A002206
and A002207 \cite{OEIS206,OEIS207}.

\section{From a generating function to a signed integral}\label{sec:integral}

\subsection{Uniform random variables and a piecewise polynomial density}
For $|x|<1$,
\begin{equation}\label{eq:uniform}
 \frac{x}{\log(1+x)}=\int_0^1(1+x)^u\,du.
\end{equation}
Let $U_1,\ldots,U_m$ be independent random variables uniform on $[0,1]$,
and put $S_m=U_1+\cdots+U_m$. Multiplying \eqref{eq:uniform} gives
\[
 F_m(x)=\E(1+x)^{S_m}.
\]
The binomial expansion is locally uniform for $S_m\in[0,m]$, so coefficient
extraction may be interchanged with the integral. Hence
\begin{equation}\label{eq:binomial_expectation}
 G_n^{(m)}=\E\binom{S_m}{n}.
\end{equation}

For $m\ge2$, the density of $S_m$ is
\begin{equation}\label{eq:density}
 f_m(t)=\frac{1}{(m-1)!}\sum_{j=0}^{m}(-1)^j\binom mj
                    (t-j)_+^{m-1},\qquad 0\le t\le m.
\end{equation}
Indeed, the unrestricted simplex
$\{u_i\ge0:\sum u_i\le t\}$ has volume $t_+^m/m!$.
Inclusion--exclusion for the conditions $u_i\le1$, followed by translation
of the selected coordinates, gives the distribution function
$\sum_j(-1)^j\binom mj(t-j)_+^m/m!$. Differentiation gives
\eqref{eq:density}.

We shall repeatedly use
\begin{equation}\label{eq:densitybounds}
 0\le f_m(t)\le\frac{t^{m-1}}{\Gamma(m)},\qquad
 f_m(t)=\frac{t^{m-1}}{\Gamma(m)}\quad(0\le t\le1),
\end{equation}
and
\begin{equation}\label{eq:density12}
 f_m(t)=\frac{t^{m-1}-m(t-1)^{m-1}}{\Gamma(m)}
                       \quad(1\le t\le2).
\end{equation}
The inequality in \eqref{eq:densitybounds} can be seen without estimating
an alternating sum: the density is an integral over a slice of the cube,
which is contained in the corresponding nonnegative simplex slice. More
explicitly, integrate over the first $m-1$ coordinates and impose
$0\le t-\sum_{i<m}u_i\le1$; dropping the upper bounds gives a region of
volume $t^{m-1}/(m-1)!$. For $m=1$, use $f_1(t)=1$ on $(0,1)$.

\subsection{The reciprocal-gamma factor}
Define
\begin{equation}\label{eq:q}
 q(t):=-\frac{1}{\Gamma(-t)}
      =\frac{\Gamma(t+1)\sin(\pi t)}{\pi}.
\end{equation}
The second equality is the gamma reflection identity
\cite[\S5.5]{DLMF}. The first expression defines an entire function;
the second is especially useful for real $t\ge0$. In particular,
\begin{equation}\label{eq:qsign}
 q(t)>0\ (0<t<1),\qquad q(t)<0\ (1<t<2),
\end{equation}
with alternating signs on subsequent unit intervals.

\begin{proposition}[Exact signed integral]\label[proposition]{prop:integral}
For every integer $n>m$,
\begin{equation}\label{eq:integral}
 a_m(n)=\int_0^m f_m(t)q(t)
                    \frac{\Gamma(n-t)}{\Gamma(n+1)}\dd.
\end{equation}
For real $x>m$, let $J_m(x)$ denote the right-hand side with $n$ replaced
by $x$. Then $J_m(n)=a_m(n)$ at the integer arguments under consideration.
\end{proposition}
\begin{proof}
The reflection formula applied to $\Gamma(t-n+1)$ yields
\[
 (-1)^{n-1}\binom tn
 =\frac{\Gamma(t+1)\sin(\pi t)}{\pi}
       \frac{\Gamma(n-t)}{\Gamma(n+1)}.
\]
The apparent exceptional values at integral $t$ follow by continuity.
Integrate this identity against $f_m$ and use
\eqref{eq:binomial_expectation}. Since $n>m$, all gamma factors in the
quotient are positive and finite on the integration interval.
\end{proof}

The distinction between $a_m$ and $J_m$ avoids assigning a real-variable
meaning to the integer sign factor $(-1)^{n-1}$.

\subsection{Existence of the threshold for every fixed order}
\begin{proposition}\label[proposition]{prop:fixedm}
For each fixed positive integer $m$,
\begin{equation}\label{eq:fixedm}
 a_m(n)=\frac{m}{n(\log n)^{m+1}}
       \left(1+O_m\!\left(\frac{1}{\log n}\right)\right)
                         \qquad(n\to\infty).
\end{equation}
In particular, $\Gr(m)$ is finite.
\end{proposition}
\begin{proof}
Uniformly on $0\le t\le m$, the gamma quotient is
$n^{-t-1}(1+O_m(n^{-1}))$. Near zero, $q(t)=t+O(t^2)$, and
\eqref{eq:densitybounds} is an equality. Fix $0<\varepsilon<1$ and put
$L=\log n$. The integral over $[0,\varepsilon]$ therefore has leading term
\[
 \frac{1}{n\Gamma(m)}\int_0^\infty t^m e^{-Lt}\dd
 =\frac{m}{nL^{m+1}}.
\]
The $O(t^2)$ term gives a relative $O_m(L^{-1})$ error. Extending the
small-interval integral to infinity costs an exponentially small quantity
in $L$ times a polynomial in $L$. On $[\varepsilon,m]$, the bounded
fixed-order amplitude contributes $O_m(n^{-1-\varepsilon})$.
This proves \eqref{eq:fixedm} and eventual strict positivity.
\end{proof}

\begin{remark}
The fixed-$m$ estimate is not uniform enough for the main question.
At the threshold, both $m$ and $\log n$ tend to infinity at comparable
rates. Substituting $m\approx\log n$ into a fixed-order asymptotic
without a new error estimate would not justify the answer.
\end{remark}

\section{A finite criterion for positivity of the whole tail}\label{sec:propagation}

Fix an integer $N>m$, and define
\begin{equation}\label{eq:prefix}
 w_{m,N}(t)=f_m(t)q(t)\frac{\Gamma(N-t)}{\Gamma(N+1)},
 \qquad F_{m,N}(u)=\int_0^u w_{m,N}(t)\dd.
\end{equation}

\begin{lemma}[Cumulative-positivity propagation]\label[lemma]{lem:propagation}
Suppose that
\begin{equation}\label{eq:prefixcondition}
 F_{m,N}(u)\ge0\quad(0\le u\le m),\qquad F_{m,N}(m)>0.
\end{equation}
Then $a_m(n)>0$ for every integer $n\ge N$.
\end{lemma}
\begin{proof}
For $n\ge N$ set
\begin{equation}\label{eq:K}
 K_{N,n}(t):=\prod_{r=N}^{n-1}\frac{r-t}{r+1}.
\end{equation}
The empty product is $1$. Gamma recurrence gives
$w_{m,n}(t)=K_{N,n}(t)w_{m,N}(t)$.
Every factor in \eqref{eq:K} is positive on $[0,m]$ and decreases with $t$.
Thus $K_{N,n}>0$ and $K'_{N,n}\le0$. Integration by parts gives
\begin{align}\label{eq:parts}
 a_m(n)
 &=K_{N,n}(m)F_{m,N}(m)
       -\int_0^m K'_{N,n}(t)F_{m,N}(t)\dd\\
 &>0.\nonumber
\end{align}
The boundary term at zero vanishes. The first displayed term is strictly
positive and the second is nonnegative, including the case $n=N$.
\end{proof}

This argument is a monotone-weight principle: a nonnegative cumulative
signed measure remains positive when reweighted toward its initial part.
Its usefulness here comes from the exact gamma recurrence, which supplies
the required decreasing weights at every later integer index.

\begin{corollary}[Integer-endpoint certificate]\label[corollary]{cor:endpoints}
It is sufficient for \eqref{eq:prefixcondition} to check
\[
 F_{m,N}(2j)\ge0\quad(1\le2j\le m),\qquad F_{m,N}(m)>0.
\]
\end{corollary}
\begin{proof}
The derivative $F'_{m,N}$ has the sign of $q$, apart from possible zeros
of the nonnegative density. Consequently $F_{m,N}$ increases on
$(0,1)$, decreases on $(1,2)$, and continues with alternating directions.
Its possible minima occur at zero, the even integer endpoints, and the
final endpoint when needed. The asserted checks therefore give
\eqref{eq:prefixcondition}.
\end{proof}

For asymptotics it is convenient to use a different sufficient condition.
\begin{corollary}[Two-interval certificate]\label[corollary]{cor:twointerval}
For $m\ge2$, suppose
\begin{equation}\label{eq:twointerval}
 F_{m,N}(2)>\int_2^m|w_{m,N}(t)|\dd.
\end{equation}
Then $a_m(n)>0$ for all integers $n\ge N$.
\end{corollary}
\begin{proof}
On $[0,1]$ the cumulative integral is nonnegative. On $[1,2]$ it decreases
to the strictly positive value $F_{m,N}(2)$. For $u\ge2$ its decrease
relative to $F_{m,N}(2)$ is bounded in magnitude by the right-hand side of
\eqref{eq:twointerval}. Thus it remains strictly positive. Apply
\cref{lem:propagation}.
\end{proof}

\section{Exact arithmetic and a finite counterexample}\label{sec:exact}

\subsection{A polynomial certificate}
For integer $N>m$, the reflection identity also gives
\begin{equation}\label{eq:polynomialw}
 w_{m,N}(t)=\frac{f_m(t)}{N!}P_N(t),\qquad
 P_N(t):=(-1)^{N-1}\prod_{r=0}^{N-1}(t-r)\in\mathbb Z[t].
\end{equation}
Thus every cumulative value at an integer endpoint is rational, without
any gamma or trigonometric evaluation. Specifically, for $1\le T\le m$,
\begin{equation}\label{eq:exactprefix}
 F_{m,N}(T)=\frac{1}{N!(m-1)!}
 \sum_{j=0}^{T-1}(-1)^j\binom mj
              \int_j^T(t-j)^{m-1}P_N(t)\dd.
\end{equation}
If $b_{j,k}$ is the coefficient of $t^k$ in
$(t-j)^{m-1}P_N(t)$ and $D=\operatorname{lcm}(1,\ldots,N+m)$, then
\begin{equation}\label{eq:integercertificate}
 \begin{split}
 N!(m-1)!D\,F_{m,N}(T)
 =\sum_{j=0}^{T-1}(-1)^j\binom mj
   \sum_{k=0}^{N+m-1}b_{j,k}\frac{D}{k+1}
             \bigl(T^{k+1}-j^{k+1}\bigr).
 \end{split}
\end{equation}
The right-hand side is an integer. Its sign is an exact certificate.

For an independent coefficient computation, write $G_n=G_n^{(1)}$.
Inverting $\log(1+x)/x$ gives
\begin{equation}\label{eq:inverse_recurrence}
 G_0=1,\qquad
 G_n=-\sum_{k=1}^n\frac{(-1)^k}{k+1}G_{n-k}.
\end{equation}
Differentiating \eqref{eq:gf} yields
$x(1+x)F'_m=m(1+x)F_m-mF_{m+1}$, hence
\begin{equation}\label{eq:order_recurrence}
 G_n^{(m+1)}=\frac{m-n}{m}G_n^{(m)}
                  +\frac{m-n+1}{m}G_{n-1}^{(m)}\quad(n\ge1).
\end{equation}
These recurrences use only rational arithmetic. They are derived and
implemented differently from the polynomial integration formula, allowing
direct cross-checks of its total integrals.

\subsection{Certified thresholds}
\begin{proposition}\label[proposition]{prop:small}
The exact thresholds include
\[
 \Gr(1)=1,\quad\Gr(2)=4,\quad\Gr(3)=11,\quad\Gr(4)=36,
 \quad\Gr(5)=113,\quad\Gr(6)=346.
\]
\end{proposition}
\begin{proof}
For $m=1$, positivity for $n\ge2$ follows directly from
\eqref{eq:integral}, since $q>0$ on $(0,1)$; also $G_1=1/2$.
For $m=2$, the recurrences give $G_3^{(2)}=0$ and
$a_2(4)=1/240$. As $F_{2,4}(2)=a_2(4)>0$, the sign pattern on the two
unit intervals and \cref{lem:propagation} give positivity for every
$n\ge4$.

For $m=3,4,5,6$, the exact certificate checks in \cref{tab:cert} establish
three facts: $a_m(B-1)<0$; every coefficient $a_m(n)$ with
$B\le n<N$ is positive; and all even-endpoint cumulative integrals at
$N$, together with the final total, are positive. The endpoint criterion
then gives positivity for every $n\ge N$. These facts imply
$\Gr(m)=B$. Each rational value is supplied in the accompanying JSON
certificate and recomputed, rather than trusted, by
\texttt{code/verify\_exact.py} using \eqref{eq:integercertificate} and
\eqref{eq:inverse_recurrence}--\eqref{eq:order_recurrence}.
\end{proof}

\begin{table}[htbp]\centering\small
\begin{tabular}{@{}rrrrll@{}}
\toprule
$m$ & $B=\Gr(m)$ & Tail anchor $N$ & Finite checks & $a_m(B-1)$ &
 Smallest checked prefix at $N$\\
\midrule
3 & 11 & 12 & 1 & $-2.74947\cdot10^{-5}$ & $2.90178\cdot10^{-5}$\\
4 & 36 & 40 & 4 & $-6.16861\cdot10^{-8}$ & $1.25186\cdot10^{-7}$\\
5 & 113 & 126 & 13 & $-1.51424\cdot10^{-9}$ & $4.63999\cdot10^{-10}$\\
6 & 346 & 378 & 32 & $-1.17700\cdot10^{-11}$ & $5.98831\cdot10^{-12}$\\
\bottomrule
\end{tabular}
\caption{Exact rational certifications. The displayed decimals are only
readable summaries of exact fractions. ``Finite checks'' counts indices
$B\le n<N$. The checked endpoints are $2,4,\ldots$ up to $m$, together
with $m$ itself.}\label{tab:cert}
\end{table}

The literal contradiction to \eqref{eq:conjecture} is therefore certified
without floating point:
\begin{equation}\label{eq:finitecounter}
 \Gr(5)=113<\frac{47}{100}3^5=\frac{11421}{100}.
\end{equation}
For this case the exact computation verifies the negative value at $112$,
the thirteen positive values at $113,\ldots,125$, and the three positive
prefixes at $2,4,5$ for $N=126$. This is a finite proof of an infinite-tail
claim. It is not a search that stops after a long observed run of positive
coefficients.

\section{The uniform large-order reduction}\label{sec:uniform}

We now let $m$ grow and place the real index in the window
\begin{equation}\label{eq:window}
 x=e^L,\qquad L=m+c,\qquad |c|\le C,
\end{equation}
where $C$ is any fixed constant. All uniform estimates in this section
refer to this compact range of $c$.

\subsection{Gamma quotient and tail bounds}
\begin{lemma}\label[lemma]{lem:gammaratio}
For sufficiently large $m$, uniformly for $0\le t\le m$ and
\eqref{eq:window},
\begin{equation}\label{eq:gammaratio}
 \frac{\Gamma(x-t)}{\Gamma(x+1)}
   =x^{-t-1}\bigl(1+O_C(m^2/x)\bigr).
\end{equation}
\end{lemma}
\begin{proof}
Write $\psi=(\log\Gamma)'$. The standard positive-real asymptotic
$\psi(y)=\log y+O(y^{-1})$ follows, for example, from
\cite[\S5.11]{DLMF}. Since $x>2m$ for all sufficiently large $m$,
\[
 \log\Gamma(x-t)-\log\Gamma(x+1)
       =-\int_{x-t}^{x+1}\psi(u)\,du.
\]
On this interval $\psi(u)=\log x+O((m+1)/x)$ uniformly.
The error after integration is $O(m^2/x)$, and exponentiation proves the
claim.
\end{proof}

For real $x>m$, define $w_{m,x}$ and $F_{m,x}$ by \eqref{eq:prefix} with
$N$ replaced by $x$.

\begin{lemma}[Exponential suppression beyond the second zero]\label[lemma]{lem:tail}
For some absolute $\eta>0$, for example $\eta=1/10$,
\begin{equation}\label{eq:tail}
 xL^m\int_2^m|w_{m,x}(t)|\dd=O_C(e^{-\eta m}).
\end{equation}
\end{lemma}
\begin{proof}
By log-convexity of $\Gamma$ on the positive axis
\cite[\S5.5]{DLMF}, for $0\le t\le m$,
\[
 \Gamma(t+1)\le\Gamma(m+1)^{t/m}\le m^t.
\]
Thus $|q(t)|\le m^t/\pi$. Combine this with
\eqref{eq:densitybounds} and \cref{lem:gammaratio}, and put
$R=L-\log m>0$. The left-hand side of \eqref{eq:tail} is bounded by a
constant times
\begin{equation}\label{eq:tail_bound_gamma}
 \frac{L^m}{\Gamma(m)}\int_2^\infty t^{m-1}e^{-Rt}\dd
 =\left(\frac LR\right)^m\Prob(X\ge2),
 \qquad X\sim\operatorname{Gamma}(m,\text{rate }R).
\end{equation}
Exponential Markov inequality with parameter $R/2$ gives
\[
 \Prob(X\ge2)\le e^{-R}\E e^{RX/2}=2^m e^{-R}.
\]
For bounded $c$ and large $m$, $(L/R)^m\le m^2$, while
$e^{-R}=m e^{-m-c}$. Therefore \eqref{eq:tail_bound_gamma} is at most
$O_C(m^3e^{-(1-\log2)m})$, which is $O_C(e^{-m/10})$.
\end{proof}

\subsection{A necessary cutoff in the gamma expectation}
Choose once and for all a smooth compactly supported real function
$\widetilde q$ which equals $q$ on $[0,2]$. Such an extension can be made
by multiplying the entire function $q$ by a smooth cutoff equal to one
on a neighborhood of $[0,2]$ and supported in $[-1,3]$.
Let
\begin{equation}\label{eq:gammaRV}
 T_{m,L}\sim\operatorname{Gamma}(m,\text{rate }L),\qquad
 p_{m,L}(t)=\frac{L^m}{\Gamma(m)}t^{m-1}e^{-Lt}\quad(t>0).
\end{equation}

\begin{proposition}[Uniform gamma-expectation reduction]\label[proposition]{prop:reduction}
Uniformly in \eqref{eq:window},
\begin{align}
 xL^mJ_m(x)&=\E\widetilde q(T_{m,L})+O_C(e^{-\eta m}),\label{eq:reductionJ}\\
 xL^mF_{m,x}(2)&=\E\widetilde q(T_{m,L})+O_C(e^{-\eta m}).\label{eq:reductionF}
\end{align}
\end{proposition}
\begin{proof}
First replace the gamma quotient on $[0,2]$ using
\cref{lem:gammaratio}. After multiplication by $xL^m$, its error is
$O_C(m^2/x)$ because $q$ is bounded on $[0,2]$ and
$f_m\le t^{m-1}/\Gamma(m)$. Next replace $f_m(t)$ there by
$t^{m-1}/\Gamma(m)$. There is no error on $[0,1]$. On $[1,2]$,
\eqref{eq:density12} bounds the absolute normalized error by
\[
 \frac{CmL^m}{\Gamma(m)}\int_1^2(t-1)^{m-1}e^{-Lt}\dd
 \le Cm e^{-L}.
\]
The substitution $u=t-1$ and extension of the $u$-integral to infinity
justify the last inequality. We have therefore obtained
\[
 xL^mF_{m,x}(2)
   =\E\bigl(q(T_{m,L})\mathbf1_{\{T_{m,L}\le2\}}\bigr)
      +O_C(m^2e^{-m}).
\]
Replacing the truncated integrand by $\widetilde q$ costs at most
$\|\widetilde q\|_\infty\Prob(T_{m,L}>2)$, which is bounded by a constant
times $2^me^{-L}$ using the same Markov argument as above. This proves
\eqref{eq:reductionF}. Adding the tail estimate \eqref{eq:tail} proves
\eqref{eq:reductionJ}.
\end{proof}

\begin{remark}[Why the cutoff cannot be omitted]\label[remark]{rem:cutoff}
The expression $\E q(T_{m,L})$ without a cutoff is not an integrable
expectation on the whole positive half-line. The amplitude
$\Gamma(t+1)\sin(\pi t)$ grows factorially, beyond the control of a fixed
exponential $e^{-Lt}$. In particular the integral of its absolute value
against the gamma density diverges. The finite original integration
range, the explicit tail estimate, and the bounded extension above are
essential. A formal replacement by an untruncated gamma expectation
would conceal a genuine analytic gap.
\end{remark}

\section{Gamma moments at a simple zero}\label{sec:moments}

The random variable in \eqref{eq:gammaRV} has mean $m/L$ and variance
$m/L^2$. In the window $L=m+c$ it is concentrated around $1$, where
$q(1)=0$ and $q'(1)=-1$. This explains why the first sign change, rather
than the later oscillations of $q$, determines the final threshold.

\subsection{Exact shifted moments and a uniform expansion}
Put $s=1/m$, and for $j\ge0$ define
\begin{equation}\label{eq:mu}
 \begin{split}
 \mu_j(s,c)
  &:=\E(T_{m,m+c}-1)^j\\
  &=\sum_{k=0}^j(-1)^{j-k}\binom jk
    \frac{\prod_{i=0}^{k-1}(1+is)}{(1+cs)^k}.
 \end{split}
\end{equation}
The empty product is $1$. This follows by expanding $(T-1)^j$ and using
$\E T^k=\rf{m}{k}/(m+c)^k$. For fixed $j$, the right-hand side is analytic
in $s$ near zero, uniformly for bounded $c$.

\begin{lemma}\label[lemma]{lem:momentsize}
For every fixed $j\ge1$,
\begin{equation}\label{eq:momentsize}
 \mu_j(s,c)=O_C(s^{\lceil j/2\rceil}).
\end{equation}
\end{lemma}
\begin{proof}
The moment-generating function of $T-1$ is
\[
 \E e^{z(T-1)}
 =\exp\left[\left(\frac m{m+c}-1\right)z
             +\sum_{r\ge2}\frac{m z^r}{r(m+c)^r}\right].
\]
The coefficient of $z$ in the exponent is $O_C(s)$, and that of $z^r$
is $O_C(s^{r-1})$ for $r\ge2$. In each product contributing to $z^j$
in the exponential, a factor of degree $1$ costs one power of $s$, while
a factor of degree $r\ge2$ costs at least $r/2$ powers. Since the total
power of $s$ is an integer, it is at least $\lceil j/2\rceil$.
Multiplying the coefficient of $z^j$ by $j!$ proves the claim.
\end{proof}

\begin{proposition}[Uniform moment expansion]\label[proposition]{prop:momentexpansion}
For every fixed $K\ge1$ there are explicit polynomials $P_r(c)$,
$1\le r\le K$, such that
\begin{equation}\label{eq:momentexpansion}
 xL^mJ_m(x)=\sum_{r=1}^{K}s^rP_r(c)+O_{C,K}(s^{K+1}),
 \qquad s=1/m,\ L=m+c,
\end{equation}
and the same expansion holds with $J_m(x)$ replaced by $F_{m,x}(2)$.
They are given by the finite coefficient formula
\begin{equation}\label{eq:Pr}
 P_r(c)=\coef{s}{r}\sum_{j=1}^{2r}
               \frac{q^{(j)}(1)}{j!}\mu_j(s,c).
\end{equation}
\end{proposition}
\begin{proof}
Taylor-expand the smooth bounded extension $\widetilde q$ at $1$ through
degree $2K+1$. Its remainder is bounded by a constant times
$|T-1|^{2K+2}$, since its derivative of order $2K+2$ is globally bounded.
The expectation of this remainder is $O_{C,K}(s^{K+1})$ by
\cref{lem:momentsize}. The signed contribution of the term of degree
$2K+1$ is also $O_{C,K}(s^{K+1})$. Thus
\[
 \E\widetilde q(T)=\sum_{j=1}^{2K}
       \frac{q^{(j)}(1)}{j!}\mu_j(s,c)+O_{C,K}(s^{K+1}).
\]
There is no constant term because $q(1)=0$. Expand the finitely many
rational functions \eqref{eq:mu} in $s$. Terms with $j>2r$ cannot
contribute to $s^r$, by \eqref{eq:momentsize}; this gives \eqref{eq:Pr}.
Finally use \cref{prop:reduction}. Its exponential errors are smaller
than every fixed power of $s$.
\end{proof}

The use of the signed odd moment in this proof is deliberate. Bounding
an odd Taylor term only by an absolute moment would generally lose a
half power. Expansion through the next odd degree and control of the
even remainder give the stated integer-order error.

\subsection{The first two moment polynomials}
The reciprocal-gamma Taylor series \cite[\S5.7]{DLMF} gives
\begin{equation}\label{eq:qseries}
 \begin{split}
 q(1+z)&=-\frac{z(1+z)}{\Gamma(1-z)}\\
 &=-z\exp\left(dz+
       \sum_{k\ge2}\frac{(-1)^{k+1}-\zeta(k)}{k}z^k\right),
                  \qquad |z|<1.
 \end{split}
\end{equation}
In particular,
\begin{align}
 q'(1)&=-1,&q''(1)&=-2d,\label{eq:qderiv12}\\
 q'''(1)&=3-3d^2+3\zeta(2),\label{eq:qderiv3}\\
 q^{(4)}(1)&=-4d^3+12d(1+\zeta(2))+8\zeta(3)-8.\label{eq:qderiv4}
\end{align}
The shifted moments needed through order $s^2$ are
\begin{align}
 \mu_1&=-cs+c^2s^2+O_C(s^3),\nonumber\\
 \mu_2&=s+(c^2-2c)s^2+O_C(s^3),\nonumber\\
 \mu_3&=(2-3c)s^2+O_C(s^3),\nonumber\\
 \mu_4&=3s^2+O_C(s^3).\label{eq:fourmoments}
\end{align}
Substitution yields
\begin{equation}\label{eq:firstexpansion}
 xL^mJ_m(x)=s(c-d)+s^2B(c)+O_C(s^3),
\end{equation}
where
\begin{equation}\label{eq:B}
 B(c)=-(1+d)c^2+2dc
       +\left(\frac13-\frac c2\right)q'''(1)
       +\frac18q^{(4)}(1).
\end{equation}
Thus $P_1(c)=c-d$ and $P_2(c)=B(c)$. A useful cancellation is
\begin{equation}\label{eq:Bd}
 B(d)=\zeta(2)+\zeta(3).
\end{equation}
This establishes both the constant shift $d$ and the first inverse-order
correction. The remaining issue is to turn this calculation into a
statement about the minimal eventual-positivity index.

\section{From the local expansion to the actual threshold}\label{sec:mainproof}

\subsection{The first correction, including the infinite tail}
We first prove
\begin{equation}\label{eq:firsttheorem}
 \log\Gr(m)=m+d-\frac{\zeta(2)+\zeta(3)}{m}+O(m^{-2}).
\end{equation}
Put $A=\zeta(2)+\zeta(3)$ and, for a fixed positive constant $M$, set
\[
 c_\pm=d-As\pm Ms^2,\qquad
 x_\pm=e^{m+c_\pm},\qquad
 N_- =\lfloor x_-\rfloor,\quad N_+=\lceil x_+\rceil.
\]
Choose the fixed compact interval $I=[d-1,d+1]$. Write
\[
 Q_s(c)=s(c-d)+s^2B(c),\qquad c_0=d-As.
\]
By \eqref{eq:Bd}, $|Q_s(c_0)|\le D_1s^3$ for some fixed $D_1$ and
all sufficiently small $s>0$. The uniform remainder in
\eqref{eq:firstexpansion}, for both the total integral and the prefix at
$2$, has absolute value at most $D_0s^3$ on $I$, for a fixed $D_0$.
Moreover $Q'_s(c)=s+s^2B'(c)\ge s/2$ on $I$ for all sufficiently small
$s$. Choose once and for all $M>2(D_0+D_1+1)$. When $s$ is small enough
that $c_0\pm Ms^2\in I$, the mean-value theorem gives
\[
 Q_s(c_+)\ge(-D_1+M/2)s^3,\qquad
 Q_s(c_-)\le(D_1-M/2)s^3.
\]
Consequently the normalized integrals at the two real indices have
opposite signs, each with an absolute margin greater than $s^3$.
The same statements hold for the normalized prefixes at $2$.

Integer rounding does not affect those signs. Indeed,
\[
 \log N_\pm-m=c_\pm+O(e^{-m})
\]
for large $m$, and the uniform expansion at these new arguments changes
the finite polynomial part by only $O(se^{-m})$. It follows that
$a_m(N_-)<0$, so
\begin{equation}\label{eq:lowerthreshold}
 \Gr(m)>N_-.
\end{equation}

For the upper bound, apply the corresponding expansion to
$F_{m,N_+}(2)$. It has the same positive normalized margin of order $s^3$.
On the other hand, \cref{lem:tail} bounds the normalized integral of
$|w_{m,N_+}|$ over $[2,m]$ by $O(e^{-\eta m})$. The positive margin
eventually dominates that tail, so \eqref{eq:twointerval} holds at $N_+$.
By \cref{cor:twointerval},
\begin{equation}\label{eq:upperthreshold}
 a_m(n)>0\quad\hbox{for every integer }n\ge N_+,
 \qquad\Gr(m)\le N_+.
\end{equation}
Combining \eqref{eq:lowerthreshold} and \eqref{eq:upperthreshold}, then
taking logarithms, proves \eqref{eq:firsttheorem}.

\begin{remark}[No assumption about global uniqueness]\label[remark]{rem:uniqueness}
This proof does not assume that $J_m(x)$ has only one zero, or that its
sign is globally monotone as a function of $x$. It finds one negative
integer coefficient below the proposed threshold and certifies the
entire tail above it. Any earlier zeros, exceptional signs, or local
oscillations are irrelevant to the resulting squeeze.
\end{remark}

\subsection{An all-orders implicit expansion}
The all-orders proof uses the same argument with a narrower bracket.
Define the formal series
\begin{equation}\label{eq:formalPhi}
 \Phi(s,c):=\sum_{r\ge1}s^{r-1}P_r(c).
\end{equation}
Its leading term is $P_1(c)=c-d$. Consequently there is a unique formal
series
\begin{equation}\label{eq:cformal}
 c(s)=d+\sum_{j\ge1}C_js^j
 \quad\hbox{such that}\quad\Phi(s,c(s))=0.
\end{equation}
There is no analytic-convergence assumption in this statement. At every
power of $s$, the next unknown $C_j$ occurs linearly with coefficient $1$,
so ordinary formal coefficient comparison determines it uniquely.
One explicit recursion is
\begin{equation}\label{eq:Crecursion}
 C_j=-\coef{s}{j}\sum_{r=2}^{j+1}s^{r-1}
        P_r\left(d+\sum_{i=1}^{j-1}C_is^i\right),\qquad j\ge1.
\end{equation}

Fix $R\ge0$ and put $c_R(s)=d+\sum_{j=1}^RC_js^j$.
Use \cref{prop:momentexpansion} with $K=R+1$. By construction the
normalized signed integral at $c=c_R(s)$ is $O_R(s^{R+2})$.
Moving $c$ to $c_R(s)\pm M_Rs^{R+1}$ changes its finite expansion by
\[
 \pm M_Rs^{R+2}(1+O_R(s)),
\]
because the derivative of its first term with respect to $c$ is $s$.
For a sufficiently large fixed $M_R$, the two normalized values have
opposite signs with polynomial-sized margins. Rounding the corresponding
indices again costs only exponentially small changes. The upper prefix
margin dominates the exponential tail, so the propagation lemma applies
exactly as before. We obtain
\[
 \log\Gr(m)=m+c_R(1/m)+O_R(m^{-R-1}),
\]
which proves \eqref{eq:allorders}.

For the first two coefficients,
\begin{equation}\label{eq:C12recursive}
 C_1=-P_2(d),\qquad C_2=-P_2'(d)C_1-P_3(d).
\end{equation}
Inserting \eqref{eq:mu}, \eqref{eq:qseries}, and \eqref{eq:Pr} gives
\[
 C_2=\zeta(2)^2+\zeta(2)\zeta(3)+\zeta(2)-\zeta(3)
                          -5\zeta(4)-3\zeta(5).
\]
Using $\zeta(2)^2=\tfrac52\zeta(4)$ gives \eqref{eq:C2}.
The supplied symbolic script reproduces these operations from the exact
finite moment formula rather than fitting coefficients to numerical data.

\subsection{Why Euler's constant occurs only in the constant shift}
Set
\[
 h(t)=e^{-d(t-1)}q(t).
\]
Near $t=1$, \eqref{eq:qseries} becomes
\begin{equation}\label{eq:hseries}
 h(1+z)=-z\exp\left(
        \sum_{k\ge2}\frac{(-1)^{k+1}-\zeta(k)}{k}z^k\right).
\end{equation}
In the local Laplace integral the factor $e^{d(t-1)}$ is absorbed by
replacing the rate $L$ with $L-d$ and multiplying by the positive constant
$e^{-d}$. More precisely, with smooth cutoffs as before,
\[
 x e^d(L-d)^mJ_m(x)
   =\E\widetilde h(T_{m,L-d})+O(e^{-\eta' m})
\]
for some $\eta'>0$ in the same compact window. This follows from the
proof of \cref{prop:reduction}; the change of normalization is bounded
above and below by positive constants there.

Write $L=m+d+u$. The formal threshold expansion in $u$ now uses an
amplitude \eqref{eq:hseries} containing no $\gamma$ or $d$ and begins at
$u=0$. Its uniqueness identifies its coefficients with the $C_j$ above.
Computing $C_j$ needs derivatives through order $2j+2$, and
\eqref{eq:hseries} shows that these belong to
$\Q[\zeta(2),\ldots,\zeta(2j+1)]$. Thus the same is true of $C_j$.
This completes the proof of every assertion in \cref{thm:main}.

\section{Further explicit consequences}\label{sec:consequences}

\subsection{The threshold itself and its successive ratios}
Exponentiation gives a convenient multiplicative form:
\begin{equation}\label{eq:multiplicative}
 \Gr(m)=e^{m+d}\left[
   1-\frac{A}{m}+\frac{C_2+A^2/2}{m^2}+O(m^{-3})\right],
 \qquad A=\zeta(2)+\zeta(3).
\end{equation}
From the all-orders result, taking one more term before differencing
also gives
\begin{equation}\label{eq:ratioexpansion}
 \log\frac{\Gr(m+1)}{\Gr(m)}
 =1+\frac{A}{m^2}+\frac{-A-2C_2}{m^3}+O(m^{-4}).
\end{equation}
To justify the last error, use \eqref{eq:allorders} through at least
$R=3$; then the individual remainders are $O(m^{-4})$. No unwarranted
extra cancellation of two uncontrolled remainders is assumed.
In particular the successive ratios are eventually greater than $e$,
although they converge to $e$. Since $e<3$, they are also eventually
less than $3$.

Numerically,
\begin{align*}
 C_1&=-2.84699097000782072187215332816\ldots,\\
 C_2&=-3.39640983572202698851434325410\ldots,\\
 C_3&=-5.67971449356146763541275716380\ldots.
\end{align*}
An exact expression for $C_3$ is recorded in \cref{sec:appendixcoeffs}.
The multiplicative approximation can be poor at small $m$; an asymptotic
expansion does not make a uniform numerical guarantee for those orders.

\subsection{Closed-form roots of the source's auxiliary polynomial}
A polynomial appearing in the source's analysis \cite{XZ2026} is
\begin{equation}\label{eq:H}
 H_m(y)=\sum_{k=1}^{\lceil m/2\rceil}(-1)^{k-1}
                         \binom m{2k-1}y^{m-2k+1}.
\end{equation}
The binomial theorem identifies this immediately as
$H_m(y)=\operatorname{Im}(y+\ii)^m$ for real $y$.
For $0<\theta<\pi$,
\begin{equation}\label{eq:cot}
 H_m(\cot\theta)=\frac{\sin(m\theta)}{\sin^m\theta}.
\end{equation}
Thus, for $m\ge2$, its complete set of roots is
\begin{equation}\label{eq:Hroots}
 \left\{\cot\frac{k\pi}{m}:k=1,\ldots,m-1\right\}.
\end{equation}
They are distinct, and their number equals its degree $m-1$.
Its largest root is therefore $\rho_m=\cot(\pi/m)$.

Consequently the associated positive real quantity is exactly
\begin{equation}\label{eq:rm}
 r_m=e^{\pi\rho_m}=\exp\!\left(\pi\cot\frac\pi m\right),
\end{equation}
and elementary Taylor expansion gives
\begin{equation}\label{eq:rmasym}
 \log r_m=m-\frac{\pi^2}{3m}-\frac{\pi^4}{45m^3}+O(m^{-5}).
\end{equation}
In particular its growth factor is also $e$. Combining this identity
with our threshold theorem yields
\begin{equation}\label{eq:comparison}
 \log\frac{\Gr(m)}{r_m}
 =d+\frac{\zeta(2)-\zeta(3)}{m}
             +\frac{C_2}{m^2}+O(m^{-3}),
 \qquad \frac{\Gr(m)}{r_m}\longrightarrow e^d.
\end{equation}
The elementary root calculation alone does not determine $\Gr(m)$;
the uniform integral analysis is what supplies the threshold constant
and the all-tail conclusion.

\section{Numerical experiments, with their limitations}\label{sec:numerical}

The file \texttt{code/numerical\_roots.py} evaluates $J_m(e^{m+c})$ using
high-precision Gauss--Legendre quadrature on each unit interval. Working
precision is increased with $m$ because direct subtraction of
$\log\Gamma(x-t)$ and $\log\Gamma(x+1)$ loses approximately $\log_{10}x$
decimal digits. The density is evaluated using its symmetry
$f_m(t)=f_m(m-t)$ to reduce cancellation.

The program finds a real zero in the predicted window and repeats the
calculation with 64 and 96 quadrature nodes per interval. It also checks
opposite numerical signs at offsets $10^{-12}$ on either side of the
computed logarithmic zero. These are consistency tests, not rigorous
interval enclosures.

Let $x_m^*$ denote the zero returned by this numerical procedure and
$c_m^*=\log x_m^*-m$. \Cref{tab:numerical} shows selected results. For
$m=3,4,5,6$, the exact certificate separately confirms that
$\lceil x_m^*\rceil$ equals the threshold. For other rows the numerical
ceiling is deliberately not labeled a certified value of $\Gr(m)$.
The theorem, rather than the quadrature, establishes the large-order
threshold law.

\begin{table}[htbp]\centering\small
\begin{tabular}{@{}rrrrr@{}}
\toprule
$m$ & $c_m^*$ & $x_m^*$ & $x_m^*/3^m$ & $x_m^*/e^m$\\
\midrule
3 & $-0.653898967$ & $10.4447664$ & $0.386843$ & $0.520014$\\
4 & $-0.439844328$ & $35.1686715$ & $0.434181$ & $0.644137$\\
5 & $-0.275925609$ & $112.626202$ & $0.463482$ & $0.758869$\\
6 & $-0.154169130$ & $345.789729$ & $0.474334$ & $0.857127$\\
7 & $-0.063651865$ & $1029.00556$ & $0.470510$ & $0.938332$\\
8 & $0.004603739$ & $2994.71318$ & $0.456442$ & $1.004614$\\
10 & $0.098337403$ & $24302.5704$ & $0.411566$ & $1.103335$\\
12 & $0.158444978$ & $190697.747$ & $0.358831$ & $1.171687$\\
20 & $0.271194889$ & $6.36309023\cdot10^8$ & $0.182492$ & $1.311531$\\
30 & $0.323892698$ & $1.47740399\cdot10^{13}$ & $0.071757$ & $1.382499$\\
50 & $0.364439498$ & $7.46445593\cdot10^{21}$ & $0.010398$ & $1.439707$\\
\bottomrule
\end{tabular}
\caption{Numerically located real zeros, not a table of certified
thresholds. The logarithmic shift is predicted to tend to
$d=0.4227843351\ldots$, and the last column to
$e^d=1.5262051116\ldots$.}\label{tab:numerical}
\end{table}

For the listed orders, the difference between the two computed values
of $c_m^*$ was less than $2\cdot10^{-54}$. That unusually close agreement
reflects a smooth integrand and high working precision; it is not a
certified error bound. The saved CSV retains additional digits and the
actual discrepancies. Both quadrature runs share the same implementation,
integral, and density formula, so their agreement is less independent than
the exact recurrence-versus-polynomial checks in \cref{sec:exact}.

The finite-order pattern explains why a base-$3$ fit can look plausible
briefly. The numerical ratio $x_m^*/3^m$ rises from order $3$ to about
order $6$ and then falls. An apparent plateau over two or three orders
cannot identify the exponential scale. The rigorous limit in
\cref{cor:base} resolves this ambiguity.

\section{Reproducibility and the trust boundary}\label{sec:reproduce}

The archive contains the manuscript in \LaTeX{} and PDF, three programs,
and their outputs. All arithmetic in the exact certificate uses Python's
standard-library integers and \texttt{fractions.Fraction}. No numerical
package is needed to verify the finite counterexample.

From the extracted archive directory, run
\begin{lstlisting}
python code/verify_exact.py
python code/derive_asymptotics.py --output data/asymptotic_coefficients.txt
python code/numerical_roots.py
pdflatex -halt-on-error gregory_sign_threshold.tex
pdflatex -halt-on-error gregory_sign_threshold.tex
\end{lstlisting}
The exact verifier recomputes the JSON data and compares every value.
The option \texttt{--write} regenerates that JSON. The symbolic script
requires SymPy, and the exploratory numerical script requires mpmath;
the versions used are recorded in \texttt{requirements.txt}. The PDF
requires a standard \TeX{} installation with the packages listed in its
preamble.

There are three distinct levels of support. The finite arithmetic checks
are exact but rely on the Python interpreter and the mathematical
propagation lemma. The symbolic checks verify the coefficient algebra
but not the analytic error estimates. The quadrature gives reproducible
numerical evidence without interval guarantees. The all-orders theorem
is proved in ordinary mathematics in
\cref{sec:uniform,sec:moments,sec:mainproof}; it does not rely on observed
floating-point signs.

No proof-assistant formalization or external peer review is claimed.
The article does not provide an explicit numerical order beyond which
the displayed asymptotic error bounds hold with specified constants.
The bounds are sufficient for the stated asymptotic theorem and eventual
contradiction. Making them effective would require carrying quantitative
constants through the cutoff derivatives and finite Taylor remainders.

\section{What is resolved, and what remains worth studying}\label{sec:outlook}

The printed base-$3$ conjecture is contradicted in two different ways.
The finite certificate resolves its literal statement. More importantly,
the large-order theorem proves a replacement growth law and rules out
any repair consisting only of changing the starting index or the two
positive constants multiplying $3^m$.

The result does not settle all questions about individual orders.
In particular, the proof does not classify all zeros of $J_m(x)$ for
$x>m$, nor all sign changes among the initial coefficients. A global
zero-counting theorem could turn the numerical continuous zero into an
exact ceiling characterization for every order, but that stronger
statement is not needed here and is not asserted.

There are several concrete continuations. One is to obtain explicit
error constants and an effective range for \eqref{eq:main}. Another is
to study late terms of the formal series and possible exponentially
small corrections, which may depend on the other unit intervals of the
uniform-sum density. A third is to extend the cumulative-measure method
to coefficients represented by other compactly supported distributions.
The key reusable property is a positive decreasing ratio between the
kernels at successive indices. Where that property holds, a local
asymptotic sign calculation can potentially certify a whole discrete
tail rather than an isolated coefficient.

\appendix
\section{The third correction and an independent coefficient recipe}\label{sec:appendixcoeffs}
For compactness, write $Z_k=\zeta(k)$. The next coefficient generated by
\eqref{eq:Crecursion} is
\begin{equation}\label{eq:C3}
\begin{split}
 C_3={}&-Z_2^3-Z_2^2Z_3-Z_2^2+9Z_2Z_3+15Z_2Z_4+6Z_2Z_5-Z_2\\
       &+6Z_3^2+5Z_3Z_4+3Z_3+4Z_4-20Z_5-35Z_6-15Z_7.
\end{split}
\end{equation}
No algebraic independence of zeta values is assumed. These are polynomial
identities in their actual real values; standard identities among even
zeta values may be used to simplify them.

Here is a direct procedure to generate as many formal corrections as
desired. Set
\[
 A_1=d,\qquad A_k=\frac{(-1)^{k+1}-Z_k}{k}\quad(k\ge2),
\]
and recursively compute
\[
 b_0=1,\qquad b_n=\frac1n\sum_{k=1}^n kA_kb_{n-k}.
\]
Then $q(1+z)=-\sum_{n\ge0}b_nz^{n+1}$, so
$q^{(j)}(1)/j!=-b_{j-1}$. Insert these values and the exact rational
moments \eqref{eq:mu} into \eqref{eq:Pr}, then apply
\eqref{eq:Crecursion}. Each desired coefficient uses finitely many
operations in a polynomial ring with rational coefficients.

For example,
\[
 P_2(c)= -\frac12\bigl(3Z_2c-3Z_2d-2Z_2-2Z_3
        +2c^2d+2c^2-3cd^2-4cd+3c+d^3+2d^2-3d\bigr).
\]
The file \texttt{data/asymptotic\_coefficients.txt} also records $P_3$ and
$P_4$ in unsimplified polynomial form. Keeping those long expressions in
a machine-readable artifact rather than expanding them throughout the
proof makes the derivation easier to audit.

\section{Certificate inventory}\label{sec:inventory}
\begin{longtable}{@{}p{0.42\linewidth}p{0.53\linewidth}@{}}
\toprule
File & Purpose\\
\midrule
\endhead
\texttt{gregory\_sign\_threshold.tex} & Complete source of this article.\\
\texttt{gregory\_sign\_threshold.pdf} & Compiled, visually inspected article.\\
\texttt{code/verify\_exact.py} & Standard-library exact rational certificates and independent total-integral cross-checks.\\
\texttt{data/exact\_certificates.json} & Full rational values for all finite checks and cumulative-prefix certificates.\\
\texttt{data/exact\_verification.log} & Output of the exact verification run.\\
\texttt{code/derive\_asymptotics.py} & Symbolic derivation from shifted gamma moments; checks the first corrections and cancellation of $d$.\\
\texttt{data/asymptotic\_coefficients.txt} & Explicit moment polynomials, correction coefficients, and decimal evaluations.\\
\texttt{code/numerical\_roots.py} & Exploratory high-precision integration with two quadrature orders.\\
\texttt{data/numerical\_roots.csv} & Real-zero estimates, normalized ratios, precision, and quadrature discrepancies.\\
\texttt{data/numerical\_run.log} & Output of the numerical run.\\
\texttt{data/thresholds\_certified.csv} & The independently certified initial threshold values, orders $1$ through $6$.\\
\texttt{README.md}, \texttt{build.sh} & Scope, instructions, and build commands.\\
\texttt{SOURCE\_AUDIT.md} & The precise target source, version, location, and search limitations.\\
\texttt{requirements.txt} & Versions used for the optional symbolic and numerical scripts.\\
\bottomrule
\end{longtable}

\clearpage
\begin{thebibliography}{9}
\bibitem{XZ2026}
Ce Xu and Jianqiang Zhao,
\emph{Rational Approximations for Reciprocals of Multiple Zeta Values and
Trivariate Cauchy Numbers}, arXiv:2609.11072v1, 10 September 2026.
Conjecture 6.1, printed p.~13; auxiliary polynomial in Lemma 4.1.
\url{https://arxiv.org/abs/2609.11072v1}.
Version checked 19 September 2026.

\bibitem{DLMF}
NIST Digital Library of Mathematical Functions,
Chapter 5, \emph{Gamma Function}, especially
\S5.5 (functional relations), \S5.7 (series expansions), and
\S5.11 (asymptotic expansions).
\url{https://dlmf.nist.gov/5.5},
\url{https://dlmf.nist.gov/5.7},
\url{https://dlmf.nist.gov/5.11}.
Accessed 19 September 2026.

\bibitem{OEIS206}
The OEIS Foundation Inc.,
\emph{The On-Line Encyclopedia of Integer Sequences}, entry A002206,
classical Gregory/Cauchy coefficient numerators.
\url{https://oeis.org/A002206}.
Accessed 19 September 2026.

\bibitem{OEIS207}
The OEIS Foundation Inc.,
\emph{The On-Line Encyclopedia of Integer Sequences}, entry A002207,
classical Gregory/Cauchy coefficient denominators.
\url{https://oeis.org/A002207}.
Accessed 19 September 2026.
\end{thebibliography}
\end{document}
